\documentclass[11pt]{article}

% Plain article style for self-study reports that do not need a conference
% layout. Use the neurips style for publication-shaped delegated reports.

\usepackage[margin=1in]{geometry}
\usepackage[T1]{fontenc}
\usepackage{amsmath,amssymb}
\usepackage{booktabs}
\usepackage{graphicx}
\usepackage[colorlinks=true,linkcolor=blue,urlcolor=blue,citecolor=blue]{hyperref}
\usepackage[numbers,sort&compress]{natbib}

\title{Why Attention Divides by the Square Root of the Key Dimension}
\author{self-study-with-ai}
\date{\today}

\begin{document}
\maketitle

\begin{abstract}
Scaled dot-product attention divides its logits by $\sqrt{d_k}$ before the
softmax~\citep{vaswani2017attention}. This report separates what the source
establishes about that choice from what it conjectures. The divisor follows
from a variance calculation: under the stated assumption of independent,
zero-mean, unit-variance components, a $d_k$-dimensional dot product has
variance $d_k$, so dividing by $\sqrt{d_k}$ restores unit variance. The
familiar explanation, that unscaled logits saturate the softmax into a region
of vanishing gradients, is offered by the source as a suspicion and is
supported by no measurement in it. This study is source-only and measures
nothing itself.
\end{abstract}

\section{Introduction}

The factor $1/\sqrt{d_k}$ appears in every implementation of transformer
attention, and the reason usually given for it is a story about softmax
saturation. This report asks a narrower question: what does the primary
source actually establish, and where does the widely repeated explanation
exceed it? The scope is one paper~\citep{vaswani2017attention} and its own
argument. Later variants and any empirical reproduction are out of scope by
the brief.

\section{Background}

Dot-product attention scores a query against each key by an inner product,
normalizes the scores with a softmax, and uses the result to weight
values. The scaled variant differs by a single constant: each logit is
divided by $\sqrt{d_k}$, where $d_k$ is the dimension of the keys
(Section~3.2.1 of~\citet{vaswani2017attention}). In multi-head attention
$d_k$ is the per-head dimension rather than the model width (Section~3.2.2),
a distinction that is easy to lose when reading the equation alone.

\section{Methods and Sources}

One source, the paper that introduced the operation, read at the published
version. Secondary explanations were excluded deliberately: they restate the
same footnote with less precision, and citing them would launder a primary
claim. The methodology is source-only, so this report contains no measurement
of its own and no experimental claim. Every statement below traces to a
locator in the single source; the anchored reading is recorded in the study's
note for that source.

\section{Findings}

\subsection{The divisor follows from a variance calculation}

The source states that if the components of a query and a key are independent
random variables with mean $0$ and variance $1$, their dot product has mean
$0$ and variance $d_k$ (Section~3.2.1, footnote~4
of~\citet{vaswani2017attention}). The variance of a sum of $d_k$ independent
terms grows linearly in $d_k$, so the standard deviation grows as
$\sqrt{d_k}$. Dividing by $\sqrt{d_k}$ therefore returns the logits to unit
variance, and no other function of $d_k$ does. This answers the primary
question: the exponent is not a tuned hyperparameter but the exact
normalization implied by the assumptions.

The result holds only under those assumptions. Learned queries and keys are
neither independent nor unit variance, so the calculation fixes a scale
rather than guaranteeing a property of a trained model.

\subsection{The motivating observation is comparative}

The source introduces the scaling in the context of a comparison. It reports
that additive attention and dot-product attention perform similarly for small
$d_k$, but that additive attention outperforms dot-product attention without
scaling for larger values of $d_k$ (Section~3.2.1). Dot-product attention is
kept because it is faster and more space-efficient in practice, being
implementable with optimized matrix multiplication.

\subsection{The gradient explanation is a conjecture in the source}

The source suspects that for large $d_k$ the dot products grow large in
magnitude, pushing the softmax into regions with extremely small gradients,
and it scales the logits to counteract that effect (Section~3.2.1). The hedge
is the source's own. No ablation isolating the scaling factor, no
gradient-magnitude measurement, and no effect size for this choice appears in
the paper. The common presentation of this mechanism as established fact is
therefore more confident than its primary source.

\section{Related Work}

Out of scope by the brief. Later proposals that revisit attention
normalization exist, and this report neither surveys nor evaluates them; a
study that did would need its own gathering pass and a wider source budget.

\section{Limitations}

This report is source-only and single-source. It measures nothing, so it
cannot say whether the gradient conjecture is correct, only that the source
marks it as a conjecture. The variance argument inherits its assumptions, and
those assumptions do not describe trained attention. Settling either open
question requires a study with an experimental methodology, which is the
natural follow-up rather than a defect of this one.

\section{Conclusion}

The $\sqrt{d_k}$ divisor is the variance normalizer for a sum of $d_k$
independent unit-variance products, under the assumptions the source
states~\citep{vaswani2017attention}. The saturation story that usually
accompanies it is a suspicion in that source and is unmeasured there. Knowing
which half is derived and which half is conjectured is the durable result of
this study.

\bibliographystyle{plainnat}
\bibliography{refs}

\end{document}
